\section{Limits of FPSs}

\subsection{Stabilization of scalars}

\begin{definition}
\label{def.fps.lim.stab}
\lean{Seq.StabilizesTo}
\leantarget
\leanok
Let $\left(  a_{i}\right)  _{i\in\mathbb{N}}=\left(
a_{0},a_{1},a_{2},\ldots\right)  \in K^{\mathbb{N}}$ be a sequence of elements
of $K$. Let $a\in K$.

We say that the sequence $\left(  a_{i}\right)  _{i\in\mathbb{N}}$
\emph{stabilizes to }$a$ if there exists some $N\in\mathbb{N}$ such that%
\[
\text{all integers }i\geq N\text{ satisfy }a_{i}=a.
\]

If $a_{i}$ stabilizes to $a$ as $i\rightarrow\infty$, then we write
$\lim\limits_{i\rightarrow\infty}a_{i}=a$ and say that $a$ is the \emph{limit}
(or \emph{eventual value}) of $\left(  a_{i}\right)  _{i\in\mathbb{N}}$.
This is legitimate, since $a$ is uniquely determined by the sequence
$\left(  a_{i}\right)  _{i\in\mathbb{N}}$.
\end{definition}

\begin{lemma}
\label{lem.fps.lim.stab-unique}
\lean{Seq.stabilizesTo_unique}
\leanhelper
\leanok
If $\left(  a_{i}\right)$ stabilizes to both $\ell_{1}$ and $\ell_{2}$,
then $\ell_{1}=\ell_{2}$.
\end{lemma}

\begin{proof}
\leanok
Let $N_{1}$ and $N_{2}$ be as in the definition for $\ell_{1}$ and $\ell_{2}$
respectively. Then $a_{\max(N_{1},N_{2})}=\ell_{1}$ and
$a_{\max(N_{1},N_{2})}=\ell_{2}$, so $\ell_{1}=\ell_{2}$.
\end{proof}

\begin{lemma}
\label{lem.fps.lim.stab-const}
\lean{Seq.stabilizesTo_const}
\leanhelper
\leanok
Any constant sequence $\left(  c,c,c,\ldots\right)  \in K^{\mathbb{N}}$
stabilizes to $c$.
\end{lemma}

\begin{proof}
\leanok
Take $N=0$. Then all $i\geq 0$ satisfy $a_{i}=c$.
\end{proof}

\begin{lemma}
\label{lem.fps.lim.stab-eventually-eq}
\lean{Seq.stabilizesTo_of_eventually_eq}
\leanhelper
\leanok
If a sequence $(a_i)$ stabilizes to $\ell$, and a sequence $(b_i)$ is
eventually equal to $(a_i)$ (i.e., there exists $N$ such that $a_i = b_i$
for all $i \geq N$), then $(b_i)$ also stabilizes to $\ell$.
\end{lemma}

\begin{proof}
\leanok
Let $N_1$ witness stabilization of $(a_i)$ to $\ell$ and $N_2$ witness the
eventual equality. Then for all $i \geq \max(N_1, N_2)$, we have
$b_i = a_i = \ell$.
\end{proof}

\begin{lemma}
\label{lem.fps.lim.stab-add}
\lean{Seq.stabilizesTo_add}
\leanhelper
\leanok
If $(a_i)$ stabilizes to $\ell_a$ and $(b_i)$ stabilizes to $\ell_b$, then
$(a_i + b_i)$ stabilizes to $\ell_a + \ell_b$.
\end{lemma}

\begin{proof}
\leanok
Let $N_a$ and $N_b$ witness the two stabilizations. For $i \geq \max(N_a, N_b)$,
we have $a_i + b_i = \ell_a + \ell_b$.
\end{proof}

\begin{lemma}
\label{lem.fps.lim.stab-mul}
\lean{Seq.stabilizesTo_mul}
\leanhelper
\leanok
If $(a_i)$ stabilizes to $\ell_a$ and $(b_i)$ stabilizes to $\ell_b$, then
$(a_i \cdot b_i)$ stabilizes to $\ell_a \cdot \ell_b$.
\end{lemma}

\begin{proof}
\leanok
Let $N_a$ and $N_b$ witness the two stabilizations. For $i \geq \max(N_a, N_b)$,
we have $a_i \cdot b_i = \ell_a \cdot \ell_b$.
\end{proof}

\begin{lemma}
\label{lem.fps.lim.stab-neg}
\lean{Seq.stabilizesTo_neg}
\leanhelper
\leanok
If $(a_i)$ stabilizes to $\ell$, then $(-a_i)$ stabilizes to $-\ell$.
\end{lemma}

\begin{proof}
\leanok
Let $N$ witness the stabilization. For $i \geq N$,
we have $-a_i = -\ell$.
\end{proof}

\begin{lemma}
\label{lem.fps.lim.stab-finset-sum}
\lean{Seq.stabilizesTo_finset_sum}
\leanhelper
\leanok
Let $S$ be a finite set, and for each $s \in S$, let $(a_{s,i})_{i \in \mathbb{N}}$
be a sequence stabilizing to $\ell_s$. Then the sequence
$\left(\sum_{s \in S} a_{s,i}\right)_{i \in \mathbb{N}}$ stabilizes to
$\sum_{s \in S} \ell_s$.
\end{lemma}

\begin{proof}
\leanok
By induction on $|S|$, using Lemma~\ref{lem.fps.lim.stab-add} for the induction step.
\end{proof}

\begin{lemma}
\label{lem.fps.lim.stab-pow-nilpotent}
\lean{Seq.stabilizesTo_pow_of_nilpotent}
\leanhelper
\leanok
If $s \in K$ satisfies $s^k = 0$ for some $k \in \mathbb{N}$ (i.e., $s$ is nilpotent),
then the sequence $(s^i)_{i \in \mathbb{N}}$ stabilizes to $0$.
\end{lemma}

\begin{proof}
\leanok
For $i \geq k$, we have $s^i = s^k \cdot s^{i-k} = 0 \cdot s^{i-k} = 0$.
\end{proof}

\subsection{Coefficientwise stabilization of FPSs}

\begin{definition}
\label{def.fps.lim.coeff-stab}
\lean{PowerSeries.CoeffStabilizesTo}
\leantarget
\leanok
Let $\left(  f_{i}\right)  _{i\in\mathbb{N}}\in
K\left[  \left[  x\right]  \right]  ^{\mathbb{N}}$ be a sequence of FPSs over
$K$. Let $f\in K\left[  \left[  x\right]  \right]  $ be an FPS.

We say that $\left(  f_{i}\right)  _{i\in\mathbb{N}}$ \emph{coefficientwise
stabilizes to }$f$ if for each $n\in\mathbb{N}$,%
\[
\text{the sequence }\left(  \left[  x^{n}\right]  f_{i}\right)  _{i\in
\mathbb{N}}\text{ stabilizes to }\left[  x^{n}\right]  f.
\]

If $f_{i}$ coefficientwise stabilizes to $f$ as $i\rightarrow\infty$, then we
write $\lim\limits_{i\rightarrow\infty}f_{i}=f$ and say that $f$ is the
\emph{limit} of $\left(  f_{i}\right)  _{i\in\mathbb{N}}$.
This is legitimate, because $f$ is
uniquely determined by the sequence $\left(  f_{i}\right)  _{i\in\mathbb{N}}$.
\end{definition}

\begin{lemma}
\label{lem.fps.lim.coeff-stab-unique}
\lean{PowerSeries.coeffStabilizesTo_unique}
\leanhelper
\leanok
If a sequence $(f_i)$ of FPSs coefficientwise stabilizes to both $\ell_1$ and
$\ell_2$, then $\ell_1 = \ell_2$.
\end{lemma}

\begin{proof}
\leanok
By extensionality (Lemma~\ref{lem.fps.ext}): for each $n$, the $n$-th coefficient of $\ell_1$ equals
that of $\ell_2$ by Lemma~\ref{lem.fps.lim.stab-unique}.
\end{proof}

\begin{lemma}
\label{lem.fps.lim.coeff-stab-const}
\lean{PowerSeries.coeffStabilizesTo_const}
\leanhelper
\leanok
A constant sequence $(g, g, g, \ldots)$ of FPSs coefficientwise stabilizes to $g$.
\end{lemma}

\begin{proof}
\leanok
For each $n$, the sequence $([x^n]g, [x^n]g, \ldots)$ is constant and stabilizes
to $[x^n]g$ by Lemma~\ref{lem.fps.lim.stab-const}.
\end{proof}

\begin{lemma}
\label{lem.fps.lim.X-pow}
\lean{PowerSeries.coeffStabilizesTo_X_pow}
\leanhelper
\leanok
The sequence $(x^i)_{i \in \mathbb{N}}$ of FPSs coefficientwise stabilizes to $0$.
\end{lemma}

\begin{proof}
\leanok
For each $n \in \mathbb{N}$, the sequence $([x^n](x^i))_{i \in \mathbb{N}}$
consists of a single $1$ (at $i = n$) and $0$s elsewhere. For $i \geq n+1$,
$[x^n](x^i) = 0$, so the sequence stabilizes to $0 = [x^n]0$.
\end{proof}

\subsection{Some properties of limits}

\subsubsection{Helpers for $x^n$-equivalence and composition}

\begin{lemma}
\label{lem.fps.lim.coeff-pow-zero}
\lean{PowerSeries.coeff_pow_eq_zero_of_constantCoeff_zero}
\leanhelper
\leanok
Let $g \in K[[x]]$ with $[x^0]g = 0$. Then for any $n, d \in \mathbb{N}$ with $d > n$,
we have $[x^n](g^d) = 0$. This follows from $\operatorname{ord}(g) \geq 1$ implying
$\operatorname{ord}(g^d) \geq d > n$.
\end{lemma}

\begin{proof}
\leanok
Since $[x^0]g = 0$, we have $\operatorname{ord}(g) \geq 1$, so
$\operatorname{ord}(g^d) \geq d > n$, and therefore $[x^n](g^d) = 0$.
\end{proof}

\begin{lemma}
\label{lem.fps.lim.coeff-subst-finite}
\lean{PowerSeries.coeff_subst_eq_finite_sum}
\leanhelper
\leanok
Let $f, g \in K[[x]]$ with $[x^0]g = 0$. Then
\[
[x^n](f \circ g) = \sum_{d=0}^{n} [x^d]f \cdot [x^n](g^d).
\]
In other words, only the first $n+1$ coefficients of $f$ contribute to $[x^n](f \circ g)$.
\end{lemma}

\begin{proof}
\leanok
For $d > n$, we have $[x^n](g^d) = 0$ by Lemma~\ref{lem.fps.lim.coeff-pow-zero},
so those terms vanish from the infinite sum.
\end{proof}

\begin{lemma}
\label{lem.fps.lim.xnEquiv-invOfUnit}
\lean{PowerSeries.xnEquiv_invOfUnit}
\leanhelper
\leanok
The relation $\overset{x^n}{\equiv}$ is compatible with taking inverses. If $f \overset{x^n}{\equiv} g$
and $u$ is a unit in $K$ such that $[x^0]f = u = [x^0]g$, then $f^{-1} \overset{x^n}{\equiv} g^{-1}$
(where $f^{-1}$ denotes the inverse of $f$ as an invertible FPS with constant term $u$).
\end{lemma}

\begin{proof}
\leanok
By strong induction on the coefficient index $k$.
For $k = 0$, both inverse coefficients equal $u^{-1}$.
For $k > 0$, the recursive formula for inverse coefficients involves only
lower-degree coefficients of the inverse (handled by the IH) and
coefficients of $f$, $g$ that agree by $x^n$-equivalence.
\end{proof}

\begin{lemma}
\label{lem.fps.lim.xnEquiv-subst}
\lean{PowerSeries.xnEquiv_subst}
\leanhelper
\leanok
The relation $\overset{x^n}{\equiv}$ is compatible with composition. If
$f_1 \overset{x^n}{\equiv} f_2$ and $g_1 \overset{x^n}{\equiv} g_2$,
with $[x^0]g_1 = [x^0]g_2 = 0$, then
$f_1 \circ g_1 \overset{x^n}{\equiv} f_2 \circ g_2$.
\end{lemma}

\begin{proof}
\leanok
By Lemma~\ref{lem.fps.lim.coeff-subst-finite}, $[x^m](f \circ g)$ depends only on
the first $m+1$ coefficients of $f$ and on powers $g^d$ for $d \leq m$.
Since $f_1 \overset{x^n}{\equiv} f_2$ gives $[x^d]f_1 = [x^d]f_2$ for $d \leq m \leq n$,
and $g_1 \overset{x^n}{\equiv} g_2$ implies $g_1^d \overset{x^n}{\equiv} g_2^d$
(by compatibility of $\overset{x^n}{\equiv}$ with powers), the result follows.
\end{proof}

\begin{theorem}
\label{thm.fps.lim.lim-crit}
\lean{PowerSeries.coeffStabilizesTo_of_forall_coeff_stabilizes}
\leantarget
\leanok
Let $\left(  f_{i}\right)  _{i\in\mathbb{N}}\in
K\left[  \left[  x\right]  \right]  ^{\mathbb{N}}$ be a sequence of FPSs.
Assume that for each $n\in\mathbb{N}$, there exists some $g_{n}\in K$ such
that the sequence $\left(  \left[  x^{n}\right]  f_{i}\right)  _{i\in
\mathbb{N}}$ stabilizes to $g_{n}$. Then, the sequence $\left(  f_{i}\right)
_{i\in\mathbb{N}}$ coefficientwise stabilizes to $\sum_{n\in\mathbb{N}}%
g_{n}x^{n}$. (That is, $\lim\limits_{i\rightarrow\infty}f_{i}=\sum
_{n\in\mathbb{N}}g_{n}x^{n}$.)
\end{theorem}

\begin{proof}
\leanok
We need to show that for each $n$, the sequence $([x^n] f_i)$ stabilizes to
$[x^n](\sum_{m \in \mathbb{N}} g_m x^m) = g_n$. But this is exactly the
hypothesis applied to $n$.
\end{proof}

\begin{lemma}
\label{lem.fps.lim.xn-equiv}
\lean{PowerSeries.exists_xnEquiv_of_coeffStabilizesTo}
\leantarget
\leanok
Assume that $\left(  f_{i}\right)  _{i\in
\mathbb{N}}$ is a sequence of FPSs, and that $f$ is an FPS such that
$\lim\limits_{i\rightarrow\infty}f_{i}=f$. Then, for each $n\in\mathbb{N}$,
there exists some integer $N\in\mathbb{N}$ such that
\[
\text{all integers }i\geq N\text{ satisfy }f_{i}\overset{x^{n}}{\equiv}f.
\]
\end{lemma}

\begin{proof}
\leanok
Let $n\in\mathbb{N}$.
For each $k\in\left\{  0,1,\ldots,n\right\}  $, we pick an $N_{k}\in
\mathbb{N}$ such that all $i\geq N_{k}$ satisfy $\left[  x^{k}\right]
f_{i}=\left[  x^{k}\right]  f$ (such an $N_{k}$ exists, since $\lim
\limits_{i\rightarrow\infty}f_{i}=f$). Then, all integers $i\geq\max\left\{
N_{0},N_{1},\ldots,N_{n}\right\}  $ satisfy $f_{i}\overset{x^{n}}{\equiv}f$.
\end{proof}

\begin{proposition}
\label{prop.fps.lim.sum-prod}
\lean{PowerSeries.coeffStabilizesTo_add, PowerSeries.coeffStabilizesTo_mul}
\leantarget
\leanok
Assume that $\left(  f_{i}\right)  _{i\in
\mathbb{N}}$ and $\left(  g_{i}\right)  _{i\in\mathbb{N}}$ are two sequences
of FPSs, and that $f$ and $g$ are two FPSs such that%
\[
\lim\limits_{i\rightarrow\infty}f_{i}=f\ \ \ \ \ \ \ \ \ \ \text{and}%
\ \ \ \ \ \ \ \ \ \ \lim\limits_{i\rightarrow\infty}g_{i}=g.
\]
Then,
\[
\lim\limits_{i\rightarrow\infty}\left(  f_{i}+g_{i}\right)
=f+g\ \ \ \ \ \ \ \ \ \ \text{and}\ \ \ \ \ \ \ \ \ \ \lim
\limits_{i\rightarrow\infty}\left(  f_{i}g_{i}\right)  =fg.
\]
\end{proposition}

\begin{proof}
\leanok
Let $n\in\mathbb{N}$.

\textbf{Addition.} The $n$-th coefficient of $f_i + g_i$ is
$[x^n]f_i + [x^n]g_i$. Since $([x^n]f_i)$ stabilizes to $[x^n]f$ and
$([x^n]g_i)$ stabilizes to $[x^n]g$, Lemma~\ref{lem.fps.lim.stab-add}
shows that $([x^n]f_i + [x^n]g_i)$ stabilizes to $[x^n]f + [x^n]g = [x^n](f+g)$.

\textbf{Multiplication.} The $n$-th coefficient of $f_i g_i$ is
$[x^n](f_i g_i) = \sum_{(a,b) \in \mathrm{antidiag}(n)} [x^a]f_i \cdot [x^b]g_i$.
Each term $[x^a]f_i \cdot [x^b]g_i$ is a product of two stabilizing sequences
(by Lemma~\ref{lem.fps.lim.stab-mul}), and a finite sum of stabilizing
sequences stabilizes (by Lemma~\ref{lem.fps.lim.stab-finset-sum}). Thus the
sequence stabilizes to $\sum_{(a,b)} [x^a]f \cdot [x^b]g = [x^n](fg)$.
\end{proof}

\begin{lemma}
\label{lem.fps.lim.neg}
\lean{PowerSeries.coeffStabilizesTo_neg}
\leanhelper
\leanok
If $(f_i)$ coefficientwise stabilizes to $f$, then $(-f_i)$ coefficientwise
stabilizes to $-f$.
\end{lemma}

\begin{proof}
\leanok
For each $n$, $([x^n](-f_i)) = (-[x^n]f_i)$ stabilizes to $-[x^n]f = [x^n](-f)$
by Lemma~\ref{lem.fps.lim.stab-neg}.
\end{proof}

\begin{lemma}
\label{lem.fps.lim.sub}
\lean{PowerSeries.coeffStabilizesTo_sub}
\leanhelper
\leanok
If $(f_i)$ coefficientwise stabilizes to $f$ and $(g_i)$ coefficientwise
stabilizes to $g$, then $(f_i - g_i)$ coefficientwise stabilizes to $f - g$.
\end{lemma}

\begin{proof}
\leanok
Write $f_i - g_i = f_i + (-g_i)$ and apply Proposition~\ref{prop.fps.lim.sum-prod}
(addition) and Lemma~\ref{lem.fps.lim.neg}.
\end{proof}

\begin{corollary}
\label{cor.fps.lim.sum-prod-k}
\lean{PowerSeries.coeffStabilizesTo_finset_sum, PowerSeries.coeffStabilizesTo_finset_prod}
\leantarget
\leanok
Let $k\in\mathbb{N}$. For each $i\in\left\{
1,2,\ldots,k\right\}  $, let $f_{i}$ be an FPS, and let $\left(
f_{i,n}\right)  _{n\in\mathbb{N}}$ be a sequence of FPSs such that%
\[
\lim\limits_{n\rightarrow\infty}\left(  f_{i,n}\right)  =f_{i}%
\]
(note that it is $n$, not $i$, that goes to $\infty$ here!). Then,%
\[
\lim\limits_{n\rightarrow\infty}\sum_{i=1}^{k}f_{i,n}=\sum_{i=1}^{k}%
f_{i}\ \ \ \ \ \ \ \ \ \ \text{and}\ \ \ \ \ \ \ \ \ \ \lim
\limits_{n\rightarrow\infty}\prod_{i=1}^{k}f_{i,n}=\prod_{i=1}^{k}f_{i}.
\]
\end{corollary}

\begin{proof}
\leanok
Follows by induction on $k$, using Proposition \ref{prop.fps.lim.sum-prod}.
\end{proof}

\begin{lemma}
\label{lem.fps.lim.isUnit-constantCoeff}
\lean{PowerSeries.isUnit_constantCoeff_of_coeffStabilizesTo}
\leanhelper
\leanok
Assume that $(g_i)$ coefficientwise stabilizes to $g$, and that each $g_i$ is
invertible (i.e., its constant coefficient is a unit). Then $g$ is also
invertible (i.e., the constant coefficient of $g$ is a unit).
\end{lemma}

\begin{proof}
\leanok
The $0$-th coefficient of $g_i$ stabilizes to the $0$-th coefficient of $g$.
For sufficiently large $i$, $[x^0]g_i = [x^0]g$. Since $[x^0]g_i$ is a unit,
so is $[x^0]g$.
\end{proof}

\begin{proposition}
\label{prop.fps.lim.sum-quot}
\lean{PowerSeries.coeffStabilizesTo_mul_inv}
\leantarget
\leanok
Assume that $\left(  f_{i}\right)  _{i\in
\mathbb{N}}$ and $\left(  g_{i}\right)  _{i\in\mathbb{N}}$ are two sequences
of FPSs, and that $f$ and $g$ are two FPSs such that%
\[
\lim\limits_{i\rightarrow\infty}f_{i}=f\ \ \ \ \ \ \ \ \ \ \text{and}%
\ \ \ \ \ \ \ \ \ \ \lim\limits_{i\rightarrow\infty}g_{i}=g.
\]
Assume that each FPS $g_{i}$ is invertible. Then, $g$ is also invertible, and
we have%
\[
\lim\limits_{i\rightarrow\infty}\dfrac{f_{i}}{g_{i}}=\dfrac{f}{g}.
\]
\end{proposition}

\begin{proof}
\leanok
The invertibility of $g$ follows from Lemma~\ref{lem.fps.lim.isUnit-constantCoeff}.
For the limit equality, we write $f_i / g_i = f_i \cdot g_i^{-1}$ and show
that the coefficients of $g_i^{-1}$ stabilize to those of $g^{-1}$.

The key step is showing that the inverse coefficients stabilize. This is done
by strong induction on the coefficient index $n$.

\textbf{Base case} ($n = 0$): The $0$-th coefficient of $g_i^{-1}$ is
$([x^0]g_i)^{-1}$. Since $[x^0]g_i$ stabilizes to $[x^0]g$, the inverse
stabilizes as well.

\textbf{Induction step} ($n > 0$): The $n$-th coefficient of $g_i^{-1}$
is determined by the recursive formula
\[
[x^n](g_i^{-1}) = -([x^0]g_i)^{-1} \sum_{\substack{(a,b) \in \mathrm{antidiag}(n) \\ b < n}}
[x^a]g_i \cdot [x^b](g_i^{-1}).
\]
By the induction hypothesis, each $[x^b](g_i^{-1})$ with $b < n$ stabilizes.
Combined with stabilization of $[x^a]g_i$ and $([x^0]g_i)^{-1}$, the
products and sums stabilize by Lemmas~\ref{lem.fps.lim.stab-mul}
and~\ref{lem.fps.lim.stab-finset-sum}.

Once we know $g_i^{-1}$ coefficientwise stabilizes to $g^{-1}$, the result
$\lim f_i g_i^{-1} = f g^{-1}$ follows from
Proposition~\ref{prop.fps.lim.sum-prod} (multiplication).
\end{proof}

\begin{proposition}
\label{prop.fps.lim.comp}
\lean{PowerSeries.coeffStabilizesTo_subst}
\leantarget
\leanok
Assume that $\left(  f_{i}\right)  _{i\in\mathbb{N}}$
and $\left(  g_{i}\right)  _{i\in\mathbb{N}}$ are two sequences of FPSs, and
that $f$ and $g$ are two FPSs such that%
\[
\lim\limits_{i\rightarrow\infty}f_{i}=f\ \ \ \ \ \ \ \ \ \ \text{and}%
\ \ \ \ \ \ \ \ \ \ \lim\limits_{i\rightarrow\infty}g_{i}=g.
\]
Assume that $\left[  x^{0}\right]  g_{i}=0$ for each $i\in\mathbb{N}$. Then,
$\left[  x^{0}\right]  g=0$ and%
\[
\lim\limits_{i\rightarrow\infty}\left(  f_{i}\circ g_{i}\right)  =f\circ g.
\]
\end{proposition}

\begin{proof}
\leanok
First, $[x^0]g = 0$ since $[x^0]g_i = 0$ for all $i$ and the $0$-th
coefficient stabilizes.

For each $n \in \mathbb{N}$, by Lemma~\ref{lem.fps.lim.xn-equiv} we can find
$N_f$ and $N_g$ such that for $i \geq \max(N_f, N_g)$, $f_i \overset{x^n}{\equiv} f$
and $g_i \overset{x^n}{\equiv} g$. By Proposition~\ref{prop.fps.xneq.comp}
(compatibility of $x^n$-equivalence with composition),
$f_i \circ g_i \overset{x^n}{\equiv} f \circ g$, so in particular
$[x^n](f_i \circ g_i) = [x^n](f \circ g)$.
\end{proof}

\begin{proposition}
\label{prop.fps.lim.deriv-lim}
\lean{PowerSeries.coeffStabilizesTo_derivativeFun}
\leantarget
\leanok
Let $\left(  f_{n}\right)  _{n\in\mathbb{N}}$ be
a sequence of FPSs, and let $f$ be an FPS such that%
\[
\lim\limits_{i\rightarrow\infty}f_{i}=f.
\]
Then,%
\[
\lim\limits_{i\rightarrow\infty}f_{i}^{\prime}=f^{\prime}.
\]
\end{proposition}

\begin{proof}
\leanok
For each $n$, $[x^n](f_i') = (n+1) \cdot [x^{n+1}]f_i$. Since $([x^{n+1}]f_i)$
stabilizes to $[x^{n+1}]f$, multiplying by the constant $(n+1)$ preserves
stabilization (by Lemma~\ref{lem.fps.lim.stab-mul} with the constant sequence
$(n+1)$). Thus the sequence stabilizes to $(n+1) \cdot [x^{n+1}]f = [x^n](f')$.
\end{proof}

\subsection{Infinite sums and products}

\begin{definition}
\label{def.fps.lim.isSummable}
\lean{PowerSeries.IsSummable}
\leanhelper
\leanok
A family $(f_n)_{n \in \mathbb{N}}$ of FPSs is \emph{summable} if for each
$n \in \mathbb{N}$, the set $\{i \in \mathbb{N} : [x^n] f_i \neq 0\}$ is finite.
\end{definition}

\begin{definition}
\label{def.fps.lim.tsum}
\lean{PowerSeries.tsum'}
\leanhelper
\leanok
The \emph{infinite sum} of a summable family $(f_n)_{n \in \mathbb{N}}$ is the FPS
$\sum_{n \in \mathbb{N}} f_n$ whose $n$-th coefficient is
$\sum_{i \in S_n} [x^n] f_i$, where $S_n = \{i : [x^n] f_i \neq 0\}$.
\end{definition}

\begin{definition}
\label{def.fps.lim.isMultipliable}
\lean{PowerSeries.IsMultipliable}
\leanhelper
\leanok
A family $(f_n)_{n \in \mathbb{N}}$ of FPSs is \emph{multipliable} if
(1) $[x^0]f_i = 1$ for all $i$, and
(2) for each $n \in \mathbb{N}$, there exists $N$ such that for all $i \geq N$
and all $k \leq n$, $[x^k]f_i = \delta_{k,0}$
(i.e., eventually all $f_i$ are $\equiv 1 \pmod{x^{n+1}}$).
\end{definition}

\begin{definition}
\label{def.fps.lim.tprod}
\lean{PowerSeries.tprod'}
\leanhelper
\leanok
The \emph{infinite product} of a multipliable family $(f_n)_{n \in \mathbb{N}}$
is the FPS $\prod_{n \in \mathbb{N}} f_n$ whose $n$-th coefficient equals
$[x^n](\prod_{j=0}^{N} f_j)$ for any sufficiently large $N$.
\end{definition}

\begin{theorem}
\label{thm.fps.lim.sum-lim}
\lean{PowerSeries.coeffStabilizesTo_partial_sum}
\leantarget
\leanok
Let $\left(  f_{n}\right)  _{n\in\mathbb{N}}$ be a
summable sequence of FPSs. Then,%
\[
\lim\limits_{i\rightarrow\infty}\sum_{n=0}^{i}f_{n}=\sum_{n\in\mathbb{N}}%
f_{n}.
\]
In other words, the infinite sum $\sum_{n\in\mathbb{N}}f_{n}$ is the limit of
the finite partial sums $\sum_{n=0}^{i}f_{n}$.
\end{theorem}

\begin{proof}
\leanok
For each coefficient index $n$, only finitely many $f_j$ have nonzero $n$-th
coefficient (by summability). Let $S$ be the finite set of such indices, and
let $N$ be the maximum element of $S$ (or $0$ if $S$ is empty). For $i \geq N+1$,
the partial sum $\sum_{j=0}^{i} [x^n]f_j$ includes all nonzero contributors,
so it equals the infinite sum $\sum_{j \in \mathbb{N}} [x^n]f_j$. Hence the
$n$-th coefficient stabilizes.
\end{proof}

\begin{theorem}
\label{thm.fps.lim.prod-lim}
\lean{PowerSeries.coeffStabilizesTo_partial_prod}
\leantarget
\leanok
Let $\left(  f_{n}\right)  _{n\in\mathbb{N}}$ be a
multipliable sequence of FPSs. Then,%
\[
\lim\limits_{i\rightarrow\infty}\prod_{n=0}^{i}f_{n}=\prod_{n\in\mathbb{N}%
}f_{n}.
\]
In other words, the infinite product $\prod_{n\in\mathbb{N}}f_{n}$ is the
limit of the finite partial products $\prod_{n=0}^{i}f_{n}$.
\end{theorem}

\begin{proof}
\leanok
By the multipliability condition, for each $n$ there exists $N$ such that
for all $j \geq N$, $f_j \equiv 1 \pmod{x^{n+1}}$ (i.e., $[x^k]f_j = \delta_{k,0}$
for $k \leq n$). Once $i$ exceeds $N$, multiplying the partial product by $f_{i+1}$
does not change the first $n+1$ coefficients, since $f_{i+1}$ acts as $1$
on those coefficients. Hence the $n$-th coefficient of the partial product
stabilizes to the $n$-th coefficient of the infinite product.
\end{proof}

\begin{corollary}
\label{cor.fps.lim.fps-as-pol}
\lean{PowerSeries.coeffStabilizesTo_trunc}
\leantarget
\leanok
Each FPS is a limit of a sequence of
polynomials. Indeed, if $a=\sum_{n\in\mathbb{N}}a_{n}x^{n}$ (with $a_{n}\in
K$), then%
\[
a=\lim\limits_{i\rightarrow\infty}\sum_{n=0}^{i}a_{n}x^{n}.
\]
More precisely, the sequence of truncations $(\operatorname{trunc}_{i+1}(a))_{i \in \mathbb{N}}$
coefficientwise stabilizes to $a$.
\end{corollary}

\begin{proof}
\leanok
For each $n$, once $i \geq n$, the truncation $\operatorname{trunc}_{i+1}(a)$
has $[x^n](\operatorname{trunc}_{i+1}(a)) = [x^n]a$ (since $n < i + 1$).
\end{proof}

\begin{theorem}
\label{thm.fps.lim.sum-lim-conv}
\lean{PowerSeries.isSummable_of_coeffStabilizesTo_partial_sum, PowerSeries.tsum'_eq_of_coeffStabilizesTo_partial_sum}
\leantarget
\leanok
Let $\left(  f_{0},f_{1},f_{2},\ldots\right)
$ be a sequence of FPSs such that $\lim\limits_{i\rightarrow\infty}\sum
_{n=0}^{i}f_{n}$ exists. Then, the family $\left(  f_{n}\right)
_{n\in\mathbb{N}}$ is summable, and satisfies%
\[
\sum_{n\in\mathbb{N}}f_{n}=\lim\limits_{i\rightarrow\infty}\sum_{n=0}^{i}%
f_{n}.
\]
\end{theorem}

\begin{proof}
\leanok
\textbf{Summability.} Fix a coefficient index $n$. Let $N$ be such that
$[x^n](\sum_{j=0}^{i} f_j) = [x^n]\ell$ for all $i \geq N$ (where $\ell$
is the limit). Suppose for contradiction that $i \geq N+1$ and $[x^n]f_i \neq 0$.
Then
\[
[x^n]\ell = [x^n]\sum_{j=0}^{i} f_j = [x^n]\sum_{j=0}^{i-1} f_j + [x^n]f_i
= [x^n]\ell + [x^n]f_i,
\]
which gives $[x^n]f_i = 0$, a contradiction. So the set
$\{i : [x^n]f_i \neq 0\} \subseteq \{0, 1, \ldots, N\}$ is finite.

\textbf{Equality.} Once summability is established, Theorem~\ref{thm.fps.lim.sum-lim}
gives $\lim_{i \to \infty} \sum_{n=0}^{i} f_n = \sum_{n \in \mathbb{N}} f_n$.
By uniqueness of limits, $\sum_{n \in \mathbb{N}} f_n = \ell$.
\end{proof}

\begin{theorem}
\label{thm.fps.lim.prod-lim-conv}
\lean{PowerSeries.isMultipliable_of_coeffStabilizesTo_partial_prod, PowerSeries.tprod'_eq_of_coeffStabilizesTo_partial_prod}
\leantarget
\leanok
Let $\left(  f_{0},f_{1},f_{2},\ldots\right)
$ be a sequence of FPSs such that $\lim\limits_{i\rightarrow\infty}\prod
_{n=0}^{i}f_{n}$ exists. Then, the family $\left(  f_{n}\right)
_{n\in\mathbb{N}}$ is multipliable, and satisfies%
\[
\prod_{n\in\mathbb{N}}f_{n}=\lim\limits_{i\rightarrow\infty}\prod_{n=0}%
^{i}f_{n}.
\]
\end{theorem}

\begin{proof}
\leanok
\textbf{Multipliability.} We first show that $[x^0]f_i = 1$ for all $i$.
Since the $0$-th coefficient of $\prod_{j=0}^{i} f_j$ stabilizes to
$[x^0]\ell$, and $[x^0](\prod_{j=0}^{i} f_j) = \prod_{j=0}^{i} [x^0]f_j$,
the product of the constant coefficients stabilizes; this forces $[x^0]f_i = 1$
for all sufficiently large $i$, and since the constant coefficients are units
whose product stabilizes, we get $[x^0]f_i = 1$ for all $i$.

Next, we need to show that for each $n$, eventually $f_i \equiv 1 \pmod{x^{n+1}}$
(i.e., $[x^k]f_i = \delta_{k,0}$ for all $k \leq n$).

We prove this by strong induction on $k$: for $0 < k \leq n$ and sufficiently
large $i$, $[x^k]f_i = 0$.

Let $P_i = \prod_{j=0}^{i-1} f_j$ denote the partial product (so that the
partial product up to $i$ is $P_i \cdot f_i$). By hypothesis, the $k$-th
coefficient of both $P_{i+1} = P_i \cdot f_i$ and $P_i$ eventually stabilizes
to $[x^k]\ell$. Hence eventually
$[x^k](P_i \cdot f_i) = [x^k]P_i$.

Expanding $[x^k](P_i \cdot f_i)$ via the convolution formula and using
$[x^0]f_i = 1$, we get
\[
[x^k]P_i + \sum_{\substack{(a,b) \in \mathrm{antidiag}(k) \\ b \neq 0, b \neq k}}
[x^a]P_i \cdot [x^b]f_i + [x^0]P_i \cdot [x^k]f_i = [x^k]P_i.
\]
Since $[x^0]P_i = 1$ (as a product of FPSs with constant term $1$), the last
term equals $[x^k]f_i$. By the induction hypothesis, the middle sum vanishes
(since $0 < b < k$ and $[x^b]f_i = 0$ for large $i$). Thus $[x^k]f_i = 0$.

Note: the formalization adds $[x^0]f_i = 1$ as an explicit hypothesis
rather than deriving it from the existence of the limit.

\textbf{Equality.} Once multipliability is established, Theorem~\ref{thm.fps.lim.prod-lim}
gives $\lim \prod_{n=0}^{i} f_n = \prod_{n \in \mathbb{N}} f_n$, and
uniqueness of limits gives the result.
\end{proof}
